% Grothendieck's question to Kostant on tori in apposition: the Borel subgroups are opposite in types A, B, C.
% Compile: tectonic -X compile article.tex
\documentclass[11pt,a4paper]{amsart}
\usepackage[T1]{fontenc}
\usepackage[utf8]{inputenc}
\usepackage{lmodern,microtype,amssymb,mathtools}
\usepackage[margin=2.7cm]{geometry}
\usepackage[hidelinks]{hyperref}

\newtheorem{theorem}{Theorem}[section]
\newtheorem{proposition}[theorem]{Proposition}
\newtheorem{lemma}[theorem]{Lemma}
\newtheorem{corollary}[theorem]{Corollary}
\theoremstyle{definition}
\newtheorem{question}[theorem]{Question}
\theoremstyle{remark}
\newtheorem{remark}[theorem]{Remark}

\newcommand{\GL}{\mathrm{GL}}
\newcommand{\PGL}{\mathrm{PGL}}
\newcommand{\C}{\mathbb C}
\newcommand{\q}[1]{\guillemotleft\,#1\,\guillemotright}
\newcommand{\leanrepo}{https://github.com/Commutator-IO/GrothendieckArchivesProofs/blob/main}

\title[Kostant's tori in apposition]{Kostant's tori in apposition:\\ the Borel subgroups are opposite in types $A$, $B$, $C$}
\author{Michel Hua}
\address{Commutator, Paris, France}
\email{michel@commutator.io}
\thanks{ORCID 0009-0002-8304-0683.}
\date{22 November 2026, revised the same day for the classical types}
\subjclass[2020]{20G15; 17B22, 22E46, 01A60}
\keywords{Kostant, cyclic element, tori in apposition, Borel subgroups, general position, Lean, Grothendieck}

\begin{document}
\emergencystretch=2em
\begin{abstract}
On 22 October 1969 Grothendieck wrote to Kostant about pairs of maximal tori $T$, $T'$ of a complex simple adjoint group that Kostant had called in apposition. Choosing Borel subgroups $B \supset T$ and $B' \supset T'$, he asked whether $B$ and $B'$ are always in general position. For $G = \PGL_n(\C)$ the answer is yes, for every $n$ and for all $n!$ Borel subgroups through $T'$: every complete flag stable under a cyclic element with nonzero coefficients is transverse to the standard flag. The proof is a short argument with a weighted shift of the coordinates, formalised in Lean~4 with Mathlib; a second proof goes through Vandermonde determinants. The same argument, with a substitute for the inverse where the cyclic element is not invertible, gives the transversality of the stable flags for $\mathrm{Sp}_{2n}$ and $\mathrm{SO}_{2n+1}$, and for $\mathrm{SO}_{2n}$ below the middle of the flag; these are formalised too. The middle step in type $D$, the exceptional types, and Grothendieck's second question on a distinguished conjugacy class of quadruples stay open. For $B_2$ and $G_2$ a computer count over finite fields found only opposite pairs.
\end{abstract}
\maketitle

\section{The letter}\label{sec:letter}

Folder~123 of the Montpellier fonds, headed \emph{Autour de Kostant}, contains on pages 21 and 22 a letter from Grothendieck to Kostant dated 22 October 1969. The letter refers to Kostant's 1959 paper on the principal three-dimensional subgroup \cite{kostant}, which attaches to a complex simple Lie algebra $\mathfrak g$ with Cartan subalgebra $\mathfrak t$ and a choice of simple roots $\alpha_1,\dots,\alpha_\ell$ the \emph{cyclic elements}
\[
E' = \sum_{i=1}^{\ell} c_i\, e_{\alpha_i} + c_0\, e_{-\theta}, \qquad c_0,\dots,c_\ell \in \C^{\times},
\]
where $\theta$ is the highest root. Such an element is regular semisimple, so its centraliser $T'$ is a maximal torus. The term \emph{in apposition} for the pair $T$, $T'$ is the letter's; we have not found it in Kostant's paper, which we could not read in full (it is behind a paywall), and we take this description of tori in apposition from the letter and from our note \cite{problems}. Let $G$ be the adjoint group, $h$ its Coxeter number, $W$ its Weyl group.

The letter asks, in the words of the page, whether for some or for all such quadruples $(T,B,T',B')$ the Borel subgroups $B$ and $B'$ are \q{in general position}, and whether there is a \emph{distinguished} conjugacy class of quadruples. It also states, without proof, that the quadruples form $|W|^2/(h\,z\,\varphi(h))$ conjugacy classes, $z$ being the order of the centre of the simply connected cover. The count follows from standard theory and is derived in our note on open problems \cite{problems}, problem \texttt{p:kostant}; we do not repeat it here.

This article answers the first question in type $A$ (Sections~\ref{sec:statement}--\ref{sec:lean}) and, in the flag form of Theorem~\ref{thm:main}, in types $B$ and $C$ and in part in type $D$ (Section~\ref{sec:classical}). Two Borel subgroups are in general position when their relative position is the longest element of $W$, that is, when $B \cap B'$ is a maximal torus. We found no statement of this for any type, but our search was partial (Section~\ref{sec:open}).

\section{Statement}\label{sec:statement}

Let $n \geq 1$ and $V = \C^n$ with basis $e_0,\dots,e_{n-1}$ and coordinates $x_0,\dots,x_{n-1}$, indices taken modulo $n$. The group is $\PGL_n(\C)$, the torus $T$ is the diagonal one and $B$ the stabiliser of the standard flag
\[
F_m = \langle e_0,\dots,e_{m-1}\rangle, \qquad 0 \leq m \leq n.
\]
The simple root vectors are $E_{i,i+1}$ for $0 \leq i \leq n-2$ and the vector of the lowest root is $E_{n-1,0}$, so a cyclic element is
\[
(E'x)_i = c_i\, x_{i+1} \qquad (0 \leq i \leq n-1,\ c_i \neq 0).
\]
Then $E'^{\,n} = (c_0\cdots c_{n-1})\cdot 1$: the eigenvalues of $E'$ are the $n$ distinct $n$-th roots of $c_0\cdots c_{n-1}$, and $E'$ is regular semisimple. The torus $T'$ is its centraliser modulo scalars. A Borel subgroup contains $T'$ exactly when it stabilises a complete flag $0 = V_0 \subset V_1 \subset \dots \subset V_n = V$ of $T'$-stable subspaces, and these are the flags of $E'$-stable subspaces; there are $n! = |W|$ of them. The flag $V_\bullet$ is opposite to $F_\bullet$ when $V_k \cap F_{n-k} = 0$ for all $k$, which is the general position of the two Borel subgroups.

\begin{theorem}\label{thm:main}
Let $c_0,\dots,c_{n-1}$ be nonzero complex numbers and $E'$ as above. For every $E'$-stable subspace $W \subseteq V$ of dimension $k$, $W \cap F_{n-k} = 0$. Consequently every Borel subgroup $B'$ containing $T'$ is in general position with $B$, and $B\cap B'$ is a maximal torus.
\end{theorem}

Up to conjugacy every quadruple of tori in apposition and Borel subgroups through them is of this form for $G=\PGL_n(\C)$, by the description of tori in apposition recalled in Section~\ref{sec:letter}. So Theorem~\ref{thm:main} says that in type $A_{n-1}$ every quadruple has $B$ and $B'$ in general position.

\section{Proof by a weighted shift}\label{sec:proof}

Define $S\colon V \to V$ by $(Sx)_j = x_{j-1}/c_{j-1}$, indices modulo $n$. It moves every coordinate up by one place, dividing by the weight.

\begin{lemma}\label{lem:inv}
$S E' = \mathrm{id}$. Hence every $E'$-stable subspace $W$ is $S$-stable.
\end{lemma}
\begin{proof}
$(SE'x)_j = (E'x)_{j-1}/c_{j-1} = c_{j-1}x_j/c_{j-1} = x_j$. The restriction of $E'$ to $W$ is injective, therefore bijective since $W$ is finite-dimensional, and $S|_W$ is its inverse.
\end{proof}

\begin{proof}[Proof of Theorem~\ref{thm:main}]
Let $W$ be $E'$-stable of dimension $k$ and let $v \in W \cap F_{n-k}$ be nonzero. Let $p$ be the largest index with $v_p \neq 0$. Since $v \in F_{n-k}$ we have $p \leq n-k-1$, hence $p + k \leq n-1$.

For $0 \leq r \leq k$ the vector $S^r v$ vanishes in all coordinates above $p+r$ and has a nonzero coordinate at $p+r$: each application of $S$ moves the support up by one place, and no index wraps around because $p+r \leq n-1$. So the vectors $v, Sv, \dots, S^k v$ have distinct highest nonzero coordinates, and are linearly independent. They lie in $W$ by Lemma~\ref{lem:inv}, and there are $k+1$ of them in a space of dimension $k$. This is a contradiction, so $v = 0$.

For the consequence, apply this to each term $V_k$ of a flag of $E'$-stable subspaces. Two Borel subgroups with opposite flags meet in a maximal torus.
\end{proof}

The proof needs only that $E'$ is invertible with a weighted shift as inverse, and it works over any field. The hypothesis on the $c_i$ is necessary: for $n=2$ and $c_1 = 0$ the only stable line is $\langle e_0\rangle = F_1$, which is not opposite to $F_1$.

\section{A second proof by Vandermonde determinants}\label{sec:vandermonde}

Let $\mu$ be an eigenvalue of $E'$, so $\mu^n = c_0\cdots c_{n-1}$. The equation $E'x = \mu x$ reads $x_{i+1} = \mu x_i / c_i$, so an eigenvector is $x_j = d_j\,\mu^j$ with $d_j = (c_0\cdots c_{j-1})^{-1}$ (and $d_0 = 1$), the weights $d_j$ not depending on $\mu$. A flag of $E'$-stable subspaces is spanned, step by step, by the eigenvectors for $\mu_{s_1},\dots,\mu_{s_n}$ in some order. The subspace $V_k$ spanned by those of a set $S$ of $k$ eigenvalues meets $F_{n-k}$ trivially exactly when the $k \times k$ matrix of the last $k$ coordinates of these eigenvectors is invertible. Its determinant is
\[
\prod_{j=n-k}^{n-1} d_j \cdot \prod_{s \in S}\mu_s^{\,n-k} \cdot \prod_{s<s' \in S}(\mu_{s'}-\mu_s),
\]
a Vandermonde determinant times nonzero factors, and it is nonzero because the $\mu_s$ are distinct.

This is the argument that comes first to mind. The note on open problems used it to check $A_2$ by hand, where all $36$ pairs are opposite. We formalised the first proof because it needs no eigenvalues.

\section{Formalisation}\label{sec:lean}

The statement is the definition \texttt{KostantOppositeTypeA}~$n$ in \href{\leanrepo/OpenProblems/AbelianVarieties.lean}{\nolinkurl{OpenProblems/\allowbreak AbelianVarieties.lean}}: for every $c$ with all $c_i \neq 0$ and every complete flag $V_\bullet$ of $E'$-stable subspaces of $\C^n$ (with $\dim V_k = k$), a vector of $V_k$ whose coordinates of index $\geq n-k$ vanish is zero. The theorem \texttt{Grothendieck.\allowbreak Folder123Kostant\allowbreak Opposite.\allowbreak kostantOppositeTypeA} proves it for every $n$ in \href{\leanrepo/Grothendieck/RepresentationTheory/Folder123KostantOpposite.lean}{\nolinkurl{Grothendieck/RepresentationTheory/Folder123\allowbreak KostantOpposite.lean}}, with Lean v4.35.0 and Mathlib; the lemmas follow Section~\ref{sec:proof} (\texttt{shiftInv\_\allowbreak cyclicElement}, \texttt{shiftInv\_\allowbreak mem}, \texttt{opposite}). The command \texttt{\#print axioms} lists only \texttt{propext}, \texttt{Classical.choice} and \texttt{Quot.sound}.

What the Lean statement leaves to this article is the translation: that the Borel subgroups through $T'$ are the stabilisers of the $E'$-stable flags, and that general position means $V_k \cap F_{n-k} = 0$. Both are standard and were argued in Section~\ref{sec:statement}, not formalised, since Mathlib has no Borel subgroups or Bruhat decomposition in this generality.

\section{The classical types}\label{sec:classical}

Let now $N = 2n$ or $2n+1$ and $V = \C^N$ with basis $e_0,\dots,e_{N-1}$, the bilinear form pairing $e_i$ with $e_{N-1-i}$ (alternating for $\mathrm{Sp}_{2n}$, symmetric for $\mathrm{SO}_{2n+1}$ and $\mathrm{SO}_{2n}$), $T$ the diagonal torus and $B$ the upper triangular Borel subgroup, which stabilises the standard flag $F_\bullet$. The cyclic element $E'$ is the sum of the root vectors of the simple roots and of the lowest root with nonzero coefficients $a_0,\dots,a_n$, written out as a matrix in the file cited below; it lies in $\mathfrak{sp}_{2n}$, $\mathfrak{so}_{2n+1}$, $\mathfrak{so}_{2n}$ respectively.

\begin{theorem}\label{thm:classical}
Let $W$ be an $E'$-stable subspace of $V$ of dimension $k$.
\begin{enumerate}
\item In type $C_n$ ($n \geq 1$) and in type $B_n$ ($n \geq 1$), $W \cap F_{N-k} = 0$.
\item In type $D_n$ ($n \geq 3$), if $W$ is isotropic and $k < n$, then $W \cap F_{2n-k} = 0$.
\end{enumerate}
\end{theorem}

In types $B$ and $C$ the Borel subgroups through $T'$ stabilise isotropic flags $V_1 \subset \dots \subset V_n$ stable under $E'$, completed by the orthogonals $V_{N-k} = V_k^{\perp}$ into complete flags of $V$ stable under $E'$; the longest element of the Weyl group acts by $-1$, and general position of two Borel subgroups is the opposition of these complete flags. So Theorem~\ref{thm:classical}~(1) gives, as Theorem~\ref{thm:main} did in type $A$, that every Borel subgroup through $T'$ is in general position with $B$; the translation is standard and is argued here, not formalised. In type $D$ the steps $k < n-1$ and the subspace $V_{n-1}$ determine the Borel subgroup, and part~(2) says they are transverse to the standard flag; we have not written out what this gives for the relative position of the two Borel subgroups.

\begin{proof}[Proof of Theorem~\ref{thm:classical}]
The proof of Theorem~\ref{thm:main} uses only an operator $S$ preserving $W$ that moves each coordinate one step along a path of coordinates, up to a nonzero factor: a vector of $W$ vanishing on the last $k$ coordinates of the path then vanishes on the whole path.

\emph{Type $C$.} In the basis $e_0,\dots,e_{2n-1}$ the matrix of $E'$ is a cyclic element of type $A_{2n-1}$, with coefficients $\pm a_i$, all nonzero, and Theorem~\ref{thm:main} applies to every $E'$-stable subspace, isotropic or not.

\emph{Type $B$.} Here $E'$ is not invertible: its kernel is the line of $u = a_0 e_0 + a_n e_{2n}$ and its image the hyperplane $\mu = 0$ for a linear form $\mu$. An explicit map $G$ satisfies $G(E'x) = x - \mu(x)\,u$; it inverts $E'$ on every $E'$-stable subspace of $\{\mu = 0\}$ and moves the coordinate $j$ to $j+1$. Let $v \in W \cap F_{N-k}$. If $E'$ is injective on $W$, then $W = E'W \subseteq \{\mu = 0\}$ is $G$-stable, and the path $0, \dots, 2n$ gives $v = 0$. Otherwise $W' = E'W$ has dimension $< k$ and contains $w = G(E'v) = v - \mu(v)\,u$; the path $0, \dots, 2n-1$ gives $w = 0$, so $v = \mu(v)\,u$, and the vanishing of the coordinate $2n$ of $v$ forces $\mu(v) = 0$.

\emph{Type $D$.} The kernel of $E'$ is spanned by $u = a_0 e_0 + a_n e_{2n-1}$ and $u' = a_{n-1} e_{n-1} - a_{n-2} e_n$, its image is $\{\mu = \mu' = 0\}$, and $G(E'x) = x - \mu(x)\,u - \mu'(x)\,u'$. The same two cases, along the paths $1, \dots, n-1, n+1, \dots, 2n-2, 0$ and $0, \dots, n-1, n+1, \dots, 2n-2$, which skip the coordinate $n$ (on the image of $E'$ it is a multiple of the coordinate $n-1$), leave $v = \mu'(v)\,u'$. There isotropy is needed: $\langle u', u' \rangle = -2\,a_{n-1}a_{n-2} \neq 0$, so $v = 0$.
\end{proof}

The statements are \texttt{KostantOppositeTypeC}, \texttt{KostantOppositeTypeB} and \texttt{KostantOppositeTypeD} in \href{\leanrepo/OpenProblems/AbelianVarieties.lean}{\nolinkurl{OpenProblems/\allowbreak AbelianVarieties.lean}}, with the three cyclic elements written as matrices, and they are proved as \texttt{kostantOppositeTypeC}, \texttt{kostantOppositeTypeB} and \texttt{kostantOppositeTypeD} in \href{\leanrepo/Grothendieck/RepresentationTheory/Folder123KostantClassical.lean}{\nolinkurl{Grothendieck/RepresentationTheory/Folder123\allowbreak KostantClassical.lean}}, where \texttt{path\_shift} is the abstract argument and \texttt{cyclicElementC\_mem\_sp}, \texttt{cyclicElementB\_mem\_so}, \texttt{cyclicElementD\_mem\_so} check that the matrices lie in the Lie algebras. As in type $A$, the translation from flags to Borel subgroups is not formalised.

In type $D$ the middle step does not hold for every maximal isotropic $V_n \supset V_{n-1}$: a numerical computation for $n = 3, 4, 5$ finds, for each of the $2^{n-1} n!$ Borel subgroups through $T'$, exactly one of the two maximal isotropic subspaces containing $V_{n-1}$ opposite to $F_n$. This computation is reported in the Lean file and was not redone for this article; the steps $k > n$, the orthogonals of the steps $k < n$, are not formalised either.

\section{What is open}\label{sec:open}

\paragraph{Other types.} For the middle step in type $D$ and for the exceptional Lie algebras, we have only a computation; the exceptional types are out of reach of the argument of Section~\ref{sec:classical}, their standard models not being matrices of a classical form. For the split adjoint groups of types $A_2$, $B_2$ and $G_2$ over $\mathbf F_q$, with $q \equiv 1 \bmod h$ and $q$ among $7,13,19,31,37$ for $A_2$ and $G_2$ and $5,13,17,29,37$ for $B_2$, the orbits of $G(\mathbf F_q)$ on quadruples number $2$, $4$ and $12$, as the formula $|W|^2/(h\,z\,\varphi(h))$ predicts, and in every orbit $\dim(\operatorname{Lie}B \cap \operatorname{Lie}B') = 2$, the rank (script \texttt{computations/kostant/kostant\_count.py} in the project repository). Since two Borel subgroups always share a maximal torus, this forces $B \cap B'$ to be a maximal torus. This checks finitely many $q$ and only the non-simply-laced types $B_2$ and $G_2$, so it proves nothing for general $q$ or for other types; for $B_2$ it agrees with Theorem~\ref{thm:classical}. We looked for the statement in the literature on cyclic elements and Coxeter tori (arXiv and Crossref searches, including Prasad and Ginzburg's papers on nearby questions) and found no source stating it, but we could not read Kostant's paper \cite{kostant}, so we do not claim the general statement is new or old.

\paragraph{The second question.} Whether some conjugacy class of quadruples is \emph{distinguished}, stable under the outer automorphisms, is not addressed. In type $A_{n-1}$ the number of classes is $((n-1)!)^2/\varphi(n)$, which is $1$, $2$, $18$ for $n = 2, 3, 4$. The letter's footnote, that $N(T)\cap N(T')$ is an extension of $(\mathbb Z/h)^{*}$ by $\mathbb Z/h \times Z(\tilde G)$, would be the starting point.

\subsection*{Declaration}
This note and the Lean file it describes were drafted with a language model (Claude, Anthropic) from the edition's registers. The author received no funding for this work.

\begin{thebibliography}{9}
\bibitem{kostant} B.~Kostant, The principal three-dimensional subgroup and the Betti numbers of a complex simple Lie group, \emph{Amer.\ J.\ Math.} 81 (1959), 973--1032.
\bibitem{springer} T.~A.~Springer, Regular elements of finite reflection groups, \emph{Invent.\ Math.} 25 (1974), 159--198.
\bibitem{problems} M.~Hua, Open questions from Grothendieck's Montpellier notes: thirteen statements a formalisation could not reach, problem \texttt{p:kostant}, 2026.
\end{thebibliography}
\end{document}
